\scanpage{157}
On peut prolonger \deleted{\unsure{bi}}multilinéairement en une application
\[
 \begin{aligned}
 E_1\times\cdots\times E_n&\mathbin{\unsure{\text{---}|\longrightarrow}} E_1\otimes\cdots\otimes E_n\\
 (x_1,\ldots x_n)&\longmapsto x_1\ast x_2\cdots\ast x_n
 \end{aligned}
\]
donné par
\[
 \mathop{\ast}_{i}x_i=\sum_{i=1}^{n}(-1)^{i\unsure{+|-}1}2^{i-1}
 \biggl(\sum\underbrace{\varepsilon_1(x_1)t_1\varepsilon_2(x_2)t_2\cdots\varepsilon_{n-i}(x_{n-i})t_{n-i}}_{\substack{n-i\ \text{facteurs}\\\varepsilon(x_j)t_j}}
 \underbrace{x_{n-i+1}\cdots x_n}_{i\ \text{facteurs }x}.
\]
Dont l'image est le sous-module de $\bigotimes E_i$ engendré par $\mathop{\ast}A_i$. Ce sous module \deleted{est} contient \inserted{donc} $\bigotimes A_i^0$ \inserted{et $t_1t_2\cdots t_n$}, et son quotient par $\bigotimes A_i^0$ est monogène, \deleted{et engendré par $t_1\cdots t_n$ (?).}
% slanted
\margin{vérifier \unsure{si un pb|ce qui se passe} en car $2$ !}

Formule plus jolie en termes des $\sigma_i$
\[
 \begin{aligned}
 x_1\ast x_2\cdots\ast x_n
 &=\sum_{\substack{i\ \text{impair}\\1\leq i\leq n}}\quad\sum_{\binom ni\text{ termes}}
 x_1\cdots x_i\,\sigma_{i+1}(x_{\unsure{i-1|i+1}})\cdots\sigma_n(x_n)
 &&(2^{n-1}\text{ termes})\\
 \sigma(x_1\ast\cdots\ast x_n)
 &=\sum_{\substack{i\ \text{pair}\\0\leq i\leq n}}\quad\sum_{\binom ni\text{ termes}}
 x_1\cdots x_i\,\sigma_{i+1}(x_{\unsure{i-1|i+1}})\cdots\sigma_n(x_n)
 &&(2^{n-1}\text{ termes})\\
 (x_1\ast\cdots\ast x_n)+\sigma(x_1\ast\cdots\ast x_n)
 &=t_1t_2\cdots t_n\,\varepsilon_1(x_1)\varepsilon_2(x_2)\cdots\varepsilon_n(x_n)
 \end{aligned}
\]
On a
\[
 (\deleted{\unsure{x_1+}}\alpha_1)\ast x_2\cdots\ast x_n=\inserted{\deleted{\unsure{x_1\ast x_2\cdots\ast x_n+}}}\alpha_1\otimes(\sigma_2(x_2)\deleted{\unsure{+}}-x_2)\otimes\cdots\otimes(\sigma_n(x_n)-x_n)
\]
si $\alpha_1\in E_1^0$ et $n-1$ formules similaires.

Si $2$ est inversible, alors $t_1\otimes\cdots\otimes t_n$ est la base d'un supplémentaire canonique de $\bigotimes A_i^0$ dans le sous-module \inserted{$E$} de $\bigotimes \deleted{\unsure{A}}E_i$ engendré par $\mathop{\ast}_i A_i$. On peut alors
% archive sheet number: 156

